\chapter{Methodology}
\label{chap:methodology}



% =======================================================================
\section{Bioverse Pipeline Overview}
\label{sec:bioverse}

% Content:
%   - Bioverse origin: Bixel & Apai 2021, extended for HZIED by Schlecker 2024.
%   - Three-stage architecture: Generator -> Survey Simulator -> Hypothesis Test.
%   - Separation of concerns: generator is expensive, survey + inference are cheap;
%     parquet-cached intermediate table sits between them.
%   - Enumerate the four extensions this thesis makes, all inside the generator:
%       (1) PLATO P1-P5 samples replace generic Gaia sphere
%       (2) FLAME ages replace uniform age prior
%       (3) four composition priors replace fixed x_H2O
%       (4) Aguichine + Hycean extend the runaway greenhouse radius model
%   - Survey simulator + hypothesis-testing inherited unchanged.


% =======================================================================
\section{The Generator}
\label{sec:generator}

% Content:
%   - Purpose: produce a "ground truth" underlying population before observational
%     selection is applied.
%   - Five substeps executed in sequence: stellar sample -> planets/orbits ->
%     runaway greenhouse radius -> stellar ages -> water-mass-fraction prior.
%   - Output: table of true planet properties handed to the survey simulator.


\subsection{Stellar sample}
\label{ssec:bioverse_stars}

% Content:
%   - Baseline: Gaia-based catalogue (Hardegree-Ullman et al. 2023) built on the
%     GCNS (Smart et al. 2021), spherical volume d_max = 120 pc.
%   - Stellar parameters (Teff, L, R, M) from Gaia + 2MASS photometry.
%   - This thesis: replace generic Gaia sphere with observed PLATO P1-P5 target
%     lists (Ch. plato_mission). Requires disabling d_max cut only.


\subsection{Planetary population and orbits}
\label{ssec:bioverse_planets}

% Content:
%   - Bergsten et al. 2022 occurrence rate priors (Kepler-derived).
%   - Per-star multiplicity, per-planet radius + orbital period drawn from prior.
%   - Semi-major axis from Kepler's third law; instellation from stellar luminosity
%     via inverse-square law.


\subsection{Runaway greenhouse radius models}
\label{sec:runaway_models}

% Content:
%   - Physical basis: Simpson-Nakajima OLR limit (Nakajima 1992, Goldblatt 2013).
%   - Baseline: Turbet 2020 mass-radius grid + Dorn-Lichtenberg 2021 water-retention
%     correction (Schlecker 2024's default).
%   - Extension 1: Aguichine 2021 for water-dominated interiors (x_H2O up to ~0.5).
%   - Extension 2: Hycean (Pierrehumbert 2022, Innes 2023) for H2-dominated
%     atmospheres. NOTE: verify the sign and value of the Hycean threshold shift.
%   - f_rgh injection: Bernoulli draw per hot-side planet decides whether the
%     inflated radius is applied.


\subsection{Stellar ages and time-since-runaway}
\label{sec:age_modelling}

% Content:
%   - Motivation: Schlecker 2024 used a uniform age prior; PLATO forecasts need
%     real per-star ages to make the atmosphere-persistence timescale meaningful.
%   - Data source: Gaia DR3 FLAME isochrone ages.
%   - Four age-assignment modes: strict, sampled, flame_flat, flat.
%   - Post-runaway age Delta t_RG = t_star - t_RG from Baraffe 1998/2015 luminosity
%     tracks.
%   - Parametric atmospheric-lifetime decay tau_atm(M_p) -> effective inflation
%     multiplier.


\subsection{Water-mass-fraction population priors}
\label{sec:composition_prior}

% Content:
%   - Motivation: replace Schlecker 2024's single injection value with a distribution.
%   - Four priors:
%       (1) dry-dominated (Rogers 2024)
%       (2) log-normal centred on Earth-like x_H2O (Simpson 2017)
%       (3) bimodal rocky/water-world (Luque & Palle 2022)
%       (4) broad log-uniform (Kimura & Ikoma 2022)
%   - Per-planet sampling of x_H2O in place of the fiducial value.


% =======================================================================
\section{Survey Simulator}
\label{sec:survey_simulator}

% Content:
%   - Purpose: turn ground-truth population into mock observed catalogue.
%   - Three steps in sequence:
%       (1) geometric transit probability p_tr = R_star / a (~0.5% for HZ around G star)
%       (2) SNR-limited detection cut based on mission photometric noise floor
%       (3) Gaussian radius-measurement noise with per-planet sigma_R
%   - Output: same column schema as generator output, fewer rows, noise-broadened
%     radius column.
%   - Both selection effects work against seeing the HZIED signature (see
%     Sec. transit_method).


% =======================================================================
\section{Hypothesis Testing}
\label{sec:inference_methodology}

% Content:
%   - Single top-level entry point: run_hzied_analysis.
%   - Two hypothesis tests + shared noise-realisation protocol described below.


\subsection{Why Bayesian model comparison}
\label{ssec:why_bayesian}

% Content:
%   - Frequentist vs Bayesian for population-level inferences.
%   - Bayesian evidence Z as marginal likelihood; ln DeltaZ as comparison statistic.
%   - Nested sampling via dynesty (Speagle 2020).
%   - Detection thresholds (Trotta 2008): moderate, strong, decisive.
%   - Schlecker 2024's headline: N >~ 100 HZ planets + f_rgh >~ 0.1 sufficient
%     for decisive detection (ln DeltaZ >~ 10).


\subsection{Full SMA-based test}
\label{ssec:sma_test}

% Content:
%   - Sliding-mean-average (SMA) construction of Schlecker 2024.
%   - Sort planets by S_obs; smooth R_obs with rolling window w = 25 neighbours.
%   - H1 free parameters: S_thresh, x_H2O, f_rgh, avg radius R_bar.
%   - H0 shares only R_bar.
%   - Priors: uniform except log-uniform x_H2O; n_live = 100.
%   - Weakness: overlapping windows -> inflated evidence; needs two-mean cross-check.


\subsection{Simple two-mean test at fixed $S_\mathrm{thresh}$}
\label{ssec:hypothesis_test_methodology}

% Content:
%   - Two-aggregate Gaussian mean-comparison (hypothesis_mode='means').
%   - Split at injected S_thresh into cold + hot subsets; each contributes one
%     aggregate (mean radius, standard error).
%   - H1: two free means (mu_cold, mu_hot); H0: single mu.
%   - Closed-form Bayes factor available when priors are broad (used in
%     Ch. kepler_tess on real data).
%   - Lower bound on SMA test; check on SMA over-confidence.


\subsection{Recovering $\hat{S}_\mathrm{thresh}$ from the posterior}
\label{ssec:sthresh_recovery}

% Content:
%   - SMA chains carry S_thresh as first parameter -> posterior median per seed.
%   - Diagnostic orthogonal to ln DeltaZ: decisive evidence + recovered threshold
%     near injection = trustworthy detection.
%   - Divergent recovered threshold = single-step model compensating for structure.
%   - Concrete example in high-f_hycean regime (ssec:hycean_mixed).


\subsection{Multi-seed noise-realisation treatment}
\label{ssec:multi_seed}

% Content:
%   - Motivation: one seed = one number is misleading (large seed-to-seed variation).
%   - Protocol: 10 noise realisations per grid point.
%   - Per-grid-point seed base s_i = s_0 + Delta*i with Delta = 200; internal
%     n_trials = 10 uses s_i + 13*k for k = 0..9.
%   - Collapse to median + 16-84% percentile band per grid point.
%   - Every Bayes factor quoted in Ch. results is a median across 10 seeds.
